\documentclass[exercice,a4]{thedocument}%
\usepackage{texmaths-tikz}%
\usepackage{pstricks}%
\usepackage{pst-bezier}%
\usepackage{pstricks-add}%
\usepackage{pst-eucl}%
\usepackage{pst-func}%
\usepackage{pst-tree}%
\usepackage{pst-circ}%
\usepackage{tikz}%
\usepackage{pgfplots}%
\usepackage{tkz-tab}%
\usepackage{tikzpagenodes}%
\usepackage{tabularray}%
\usepackage{tabularx}%
\usepackage{multirow}%
\usepackage{multicol}%
\usepackage{scratch3}%
\usepackage{enumitem}%
\usepackage[np]{numprint}%
\usepackage{pas-tableur}%
\usepackage{listings}%
\usepackage{graphicx}%
\usepackage{adjustbox}%
\usepackage{needspace}%
\usepackage{xcolor}%
\usepackage{eurosym}%
\usepackage{diagbox}%
\usepackage{colortbl}%
\usepackage{multido}%
\usepackage{icomma}%
\usepackage[french]{babel}%
\graphicspath{{../brevet/2026/}}%
\startdatakeys
    \source{26GENMATAG1}
    \BO{2026}
    \cycle{4}
    \competencesDuSocle{CA,RA,CO}
    \objectifsApprentissage{C4A1aa,C4A2ad,C4B2cb,C4D2aa,C4D3ac}
    \motsCles{fonction linéaire,fonction affine,graphique}
\enddatakeys
\definecolor{gris}{gray}{0.5}%
\providecommand{\wideparen}[1]{\overset{\frown}{#1}}%
\begin{document}%
\providecommand{\degres}{^\circ}%
Un centre de loisirs organise pour les jeunes d'une commune des activités pendant les vacances.

Pour s'inscrire dans ce centre de loisirs, plusieurs tarifs sont proposés :
\begin{itemize}
	\item Tarif A : on paie 5,50~\euro{} par demi-journée passée dans le centre de loisirs ;
	\item Tarif B : on paie un abonnement annuel de 15~\euro{} puis 3,50~\euro{} pour chaque demi-journée passée dans le centre ;
	\item Tarif C : on paie un forfait annuel de 64~\euro{} donnant accès à autant de demi-journées que l'on veut.
\end{itemize}

\begin{enumerate}
	\item Quel est le prix pour 12 demi-journées avec le tarif B ?
	\item On désigne par $x$ le nombre de demi-journées passées dans le centre de loisirs.
	On considère les trois fonctions $f,\, g$, et $h$ suivantes :

	\[f : x \longmapsto 15 + 3,5x \qquad\qquad g : x \longmapsto 64 \qquad\qquad h : x \longmapsto 5,5x\]

	Sans justifier, associer chacune de ces fonctions au tarif qu'elle représente.

	\item Le graphique ci-dessous représente les fonctions $f,\, g$, et $h$.
	Sans justifier, associer à chaque représentation graphique ($d_1$), ($d_2$) et ($d_3$) la fonction $f$, $g$ ou $h$ correspondante.

	\begin{center}
		%Début Tikz
		\begin{tikzpicture}[x=11mm, y=1.1mm, >=stealth]
		\draw[color = gray!35, line width = 0.3pt, xstep=0.2, ystep = 2] (0,0) grid (11.4,84);
		\draw[color = gray!55, line width = 0.5pt, xstep=1, ystep = 10] (0,0) grid (11.4,84);
		\draw[->,line width=1pt] (0,0) -- (11.4,0) node[above left,text width=2.4cm,align = center] {\footnotesize Nombre de demi-journées};
		\draw[->,line width=1pt] (0,0) -- (0,84) node[below right] {Prix en \euro{}};
		\foreach \x in {1,...,11}
			\draw (\x,2pt) -- (\x,-2pt) node[below] {\np{\x}};
		\foreach \y in {10,20,...,80}
			\draw (2pt,\y) -- (-2pt,\y) node[left] {\np{\y}};
		\draw[teal,line width=1.4pt] (0,64) -- (11.4,64) node[pos=0.4,above]{$(d_2)$};
		\draw[blue,line width=1.4pt] (0,0) -- (11.4,62.7) node[pos=0.4,below]{$(d_3)$};
		\draw[red,line width=1.4pt] (0,15) -- (11.4,54.9)
		node[pos=0.4,above]{$(d_1)$};
		\end{tikzpicture}
		%Fin Tikz
%		%Début PSTricks
%		\psset{xunit=1cm,yunit=0.1cm,arrowsize=2pt 3}
%		\begin{pspicture}(-0.6,-5)(11,85)
%			\psaxes[linewidth=1.25pt,Dy=10,labelFontSize=\scriptstyle]{->}(0,0)(0,0)(11,85)
%			\multido{\n=0+1}{12}{\psline[linewidth=0.15pt](\n,0)(\n,85)}
%			\multido{\n=0+10}{9}{\psline[linewidth=0.15pt](0,\n)(11,\n)}
%			\psplot[plotpoints=2000,linewidth=1.25pt,linecolor=red]{0}{11}{3.5 x mul 15 add}
%			\psplot[plotpoints=2000,linewidth=1.25pt,linecolor=blue]{0}{11}{5.5 x mul}
%			\psline[linewidth=1.25pt,linecolor=green](0,64)(11,64)
%			\uput[u](9,0){\footnotesize Nombre de demi-journées}
%			\uput[r](0,83){\footnotesize Prix en \euro}
%			\rput(3,14){\blue $(d_3)$}\rput(3,30){\red $(d_1)$}\rput(3,66){\green $(d_2)$}
%		\end{pspicture}
		%Fin PSTricks
	\end{center}

	\item La droite ($d_1$) représente-t-elle une situation de proportionnalité ?

	\textit{Détailler la réponse en précisant la démarche.}

	\item Lire graphiquement le nombre de demi-journées à partir duquel le tarif B devient plus avantageux que le tarif A. Aucune justification n'est demandée.

	\needspace{10\lineskip}
	\item On a réalisé dans un tableur le tableau de valeurs correspondant au tarif B suivant le nombre de demi-journées de présence.

	\begin{tblr}{hlines, vlines, colsep=1pt,
			colspec={*{18}{X}},
			columns={c,font=\sffamily},	rows={m},
			cell{2-Z}{2-Z}={font=\small \sffamily},
			row{1}={rowsep=0pt,ht=7mm, bg=gray9},
			column{1}={colsep=0pt,7mm,bg=gray9},
			column{2}={2.8cm},
			cell{1}{1}={gray9,halign=r,valign=b,cmd={\tikz[baseline={(0,0.15)}]{\clip(0,0) rectangle (0.6,0.5) ;\fill[gray!50] (0,0)--(0.5,0.5)--(0.5,0);}}}
			}
		  & A                                   & B    & C  & D    & E  & F    & G  & H    & I  & J    & K  & L    & M  & N    & O  & P    & Q   \\
		1 & Nombre de demi-journées d'activités & 1    & 2  & 3    & 4  & 5    & 6  & 7    & 8  & 9    & 10 & 11   & 12 & 13   & 14 & 15   & 16  \\
		2 & Tarif B en \euro                    & 18,5 & 22 & 25,5 & 29 & 32,5 & 36 & 39,5 & 43 & 46,5 & 50 & 53,5 & 57 & 60,5 & 64 & 67,5 & 71
	\end{tblr}

	\begin{enumerate}
		\item Parmi les trois formules proposées ci-dessous, recopier celle qui a été saisie dans la cellule \textsf{B2}, puis étirée vers la droite pour compléter le tableau.

		\begin{center}
			\begin{tblr}{colspec={*{3}{X[c,font=\sffamily]}},hlines, vlines}
				=15+3,5*$x$ & =15+3,5*1 & =15+3,5*B1 \\
			\end{tblr}
		\end{center}

		\item En utilisant le tableau de valeurs ci-dessus, déterminer le nombre de demi-journées à partir duquel le tarif C devient plus avantageux que le tarif B. Aucune justification n'est demandée.
	\end{enumerate}
\end{enumerate}

%---------------------------------------------------------
% Fin du document
%---------------------------------------------------------
\medskip
\showthesource
\end{document}
